\addcontentsline{toc}{chapter}{Requirements}
\section*{\Huge Requirements}
\markright{Requirements}

\section{Intended Users}
% In this section, you identify the stakeholders of the project. Who will be using the system, and how will they use it? Are there different groups of users requiring different things for their respective use cases? If so, label these groups for later referral and explain the differences. For research projects (ML, classifier, etc.), the use cases can be abstract (for example: speeding up the data processing time for ML researchers). If the project is sponsored, the use case and system functionality sections must explain how these requirements align with the company's needs.

\begin{itemize}
    \item \textbf{Machine Learning Researchers and Benchmark Evaluators:}
Researchers working on text-to-SQL systems who need a reproducible, modular framework to test hypotheses about schema retrieval, query planning, and self-correction strategies. Their use case will be running the pipeline against Spider 2.0 Lite \cite{lei2024spider} with different configurations (e.g., enabling/disabling query planning) and comparing Execution Accuracy across ablations. Their primary requirement is a well-documented, configurable code base with clear experiment tracking.

\item \textbf{Data and Analytics Teams: }
Non-technical business analysts and technical stakeholders like data engineers in enterprises who need to extract insights from SQL databases without writing queries manually. Their use case is a pipeline designed to handle complex, real-world schemas which generally comprise of large column counts, nested types, multiple dialects that these users encounter. A system that works only on simple schemas would not serve this audience. Their primary requirement is an accurate and robust system that works well on enterprise level complex schemas.
\end{itemize}

\section{Functional Requirements}

Our system provides a framework to convert natural-language questions into executable SQL over complex enterprise-style schemas (including nested types and multiple SQL dialects), evaluate outputs, and support reproducible experimentation for research and benchmarking.

\begin{itemize}
    \item \textbf{Natural Language Query Intake}: The system shall accept user questions through CLI/batch mode (benchmark runs) and validate input format, schema target, and selected dialect. Alternate flows include if the question is empty or invalid, the system returns an actionable validation feedback. If the schema/database is unavailable, the system returns an error and logs incidents.
    \item \textbf{Schema Understanding and Retrieval}: The system shall retrieve relevant tables or columns or relations from very large schemas and support semantic retrieval for nested and ambiguous schema elements. The system shall provide a ranked list of retrieved schema items with relevance scores (or confidence indicators). Alternate flows include if the schema context exceeds token budget, the system performs context compression and summarization and records truncation events.
    \item \textbf{Query Planning}: The system shall support an explicit query planning stage that can be enabled or disabled via configuration. When enabled, the system shall generate a structured intermediate plan (sub-goals, joins, filters, aggregation strategy) before SQL synthesis. The system shall persist plan artifacts for later analysis. Alternate flows include if the planner output is malformed or incomplete, the system would retry the planner with a constrained format prompt. If the planner repeatedly fails, the system would bypass planning and proceed with direct generation, logging fallback.
    \item \textbf{SQL Generation}: The system shall generate dialect-aware SQL (e.g., BigQuery, Snowflake, BigQuery). The system shall enforce output constraints (valid SQL block, no destructive commands unless explicitly authorized) and shall attach provenance metadata: model, prompt version, config ID, timestamp. Alternate flows include if generated SQL is syntactically invalid, system invokes correction stage. If SQL references non-existent schema items, the system loops back to schema retrieval/planning.
    \item \textbf{SQL Validation and Execution}: The system shall run syntax and schema checks prior to execution, and execute SQL on target database or benchmark execution environment. The system shall capture execution status, runtime, and result metadata. Alternate flows include on execution error, the system would classify error type (syntax, semantic, permission, timeout, dialect mismatch). On timeout or resource failure, the system would apply retry policy. 
    \item \textbf{Iterative Self-Correction Loop}: The system shall support multi-turn corrective refinement of SQL using execution feedback. The number of correction iterations, stopping criteria, and retry strategies shall be configurable. Each correction iteration shall be logged with prior SQL, error message, revised SQL, and rationale trace. Alternate flows are if corrections fail after max iterations, the system returns the best attempt SQL with failure classification. If correction converges to executable SQL, the system marks iteration trace as successful and stores final output.
    \item \textbf{Error Analysis and Diagnostics}: The system shall produce structured error reports (CSV/JSON) for incorrect queries. The logs would have metrics for aggregate complexity comparison, CTE statistics (and small CTE example extracts), metric distributions, correlations with correctness, structural risk patterns (e.g. high complexity, nesting depth, join count, length, readability thresholds) as percentages and risk ratios. Alternate flows are if logs are incomplete, system flags the run as non-comparable for research reporting.
    \item \textbf{Configuration and Modularity}: The system shall expose a configuration system for model provider selection, stage enable or disable flags, prompt templates or versions, retrieval strategy, correction limits, evaluation settings. The architecture shall allow replacing components like retriever, planner, generator, corrector independently. Alternate flows are, upon invalid configuration, the system performs schema validation and reports exact invalid fields before run start.
    \item \textbf{Logging, Auditability, and Reproducibility}: The system shall log every pipeline stage with deterministic run IDs and capture environment metadata (code version, dataset version, model version, configuration hash). The system shall support reproducing a historical run from stored artifacts.
\end{itemize}

\section{Non-Functional Requirements} \label{sec:nfr}
%How and in what manner will the mentioned system functionalities be carried out? This section lays out expectations and/or constraints on speed, reliability, security, etc. (i.e., requirements pertaining to the system's characteristics that do not directly pertain to the user).

%Examples of non-functional requirements are tolerance, time (speed) constraints, the size of data, transparency and security considerations, etc. Whenever possible, mention use cases or scenarios you are thinking of when defining your non-functional requirements. Additionally, make sure to include the definition of the non-functional requirements you are using in the Definitions section below.

\begin{itemize}
    \item \textbf{Latency}: The end-to-end workflow from natural language input to SQL output is expected to complete in 30 seconds. Each workflow may consist of multiple iterations, including LLM call, error feedback, SQL validation, which should take between 5 to 8 seconds. The RAG retrieval for selecting tables and columns is expected to complete within 2 seconds. These latency constraints will be managed by a timeout control. If each iteration exceeds the time constraint or reaches the max iteration, the system should now return a result with an indicator of low confidence. This ensures that when user enter their input question, the system is able to return a validated SQL output in a reasonable time manner regardless of the number of self-correction iterations.
    \item \textbf{Scalability}: To meet the requirements of the enterprise-scale schemas like in the Spider 2.0-Lite benchmark, the system is expected to handle database schema of up to 1,000 columns per database, with tolerance of up to 3,000 columns. It should also be capable of processing multi-dialect databases, such as Snowflake, BigQuery, MySQL, as well as the given metadata and documentation provided by the benchmark. Upon the database connection, the system must build RAG index construction within 10 minutes for schemas up to 1,000 columns as a one-time preprocessing step.
    \item \textbf{Execution Accuracy}: The current top five execution accuracy of Spider2.0-lite range from 55.21 to 69.65. Since the goal is to be in the top five of this benchmark, we aim to achieve a score greater than 55 with the current framework as the benchmark at 43. The self-correction loop also present non-decreasing accuracy, as decreasing accuracy only increases latency for users.
    \item \textbf{Transparency}: Each user input should produce a trace log that includes retrieved context from RAG, generated SQL, any execution or validation errors, and correction reasoning. Proper and detailed logs can help improve debugging and understand how each step leads to an improvement or a bottleneck.
    \item \textbf{Security}: The system should never execute write queries like INSERT, UPDATE, DROP, and DELETE against the target database. Every generated SQL should be validated to ensure that it is a read-only execution. For example, if the user attempts to delete some data, the system should return an appropriate error message instead of going through the SQL generation and execution.
\end{itemize}

\section{Resource Requirements}
%What data is needed to undertake the project? What hardware is needed? What non-development human effort is needed (e.g., manual text annotation, testing, etc) What cloud resources are needed for this project, and what is the cost? (e.g., rates, GPU/CPU hours, etc.) This section should mention all the resources clearly necessary, given the items discussed in the rest of the document.
\begin{itemize}
    \item \textbf{Data}: The project uses the Spider 2.0-lite benchmark with pre-existing SQLite schema files, gold SQL, external knowledge, and LLM-generated data/database summaries already in our NL2SQL repository. The only other data involved will be intermediate analysis artifacts like a one-time offline PK/FK relationship extraction across the databasees, and possibly a vector DB of all the schema embedded into a semantic space.
    \item \textbf{Hardware}: The project follows a three-phase model strategy, and hardware requirements evolve with it. In Phase 1 (local development and rapid iteration), SQL generation runs on Qwen3-8B served via vLLM on a single NVIDIA A6000 (48GB VRAM) GPU node on PSC. Schema enrichment is a one-time offline job using Qwen3-32B-AWQ (4-bit quantised), which fits comfortably within the same node at around 20GB VRAM. Embeddings are produced by nomic-embed-text-v1.5 on CPU or a lightweight GPU process, and the resulting vector index (\texttt{dce.duckdb}) is a single $\sim$500MB file covering all 29 databases in the SQLite track. Small validation runs of 5--10 questions can be completed on any single GPU node in under 15 minutes; full 135-question benchmark runs require the vLLM server running continuously for approximately 3 hours. In Phase 2, we scale to larger open-source models (Qwen3-72B-AWQ or similar), which may require multi-GPU nodes or higher-memory instances on PSC. In Phase 3, SQL generation moves entirely to frontier model APIs (Gemini Pro or GPT-4o), which eliminates the local GPU inference requirement. At that point only the lightweight embedding server and the pre-built vector index remain on-device.
    \item \textbf{Non-Development Human Effort}: No explicity data annotation or crowdsourcing required. There will be a degree of iterative manual analysis (along with detailed error analysis dashboards) after benchmark runs for prompt and framework improvements, but nothing otherwise.
    \item \textbf{Cloud Resources \& Cost}: Token estimates for a single full benchmark run (135 SQLite instances) are $\sim$7.5M input tokens and $\sim$1.2M output tokens. Based on this, and current API pricings, the overall costs are captured in the below tables.
    \begin{table}[H]
\centering
\begin{tabular}{|l|r|r|r|}
\hline
\textbf{Resource} & \textbf{Cost per Run} & \textbf{Expected Runs} & \textbf{Total Estimated Cost} \\
\hline
Cloud VM (e2-standard-4) & \$2 & 45 & \$90 \\
\hline
Gemini 3.1 Flash (functional validation) & \$4 & 15 & \$60 \\
\hline
Gemini 3.1 Pro (full benchmark) & \$30 & 15 & \$450 \\
\hline
GPT-5.4 (full benchmark) & \$40 & 10 & \$400 \\
\hline
\textbf{Total} & & & \textbf{\$1000} \\
\hline
\end{tabular}
\caption{Estimated cloud resource costs for the project.}
\label{tab:cloud-costs}
\end{table}

\begin{table}[H]
\centering
\begin{tabular}{|l|r|r|}
\hline
\textbf{Model} & \textbf{Input (per 1M tokens)} & \textbf{Output (per 1M tokens)} \\
\hline
Gemini 3.1 Flash & \$0.25 & \$1.50 \\
\hline
Gemini 3.1 Pro & \$2.00 & \$12.00 \\
\hline
GPT-5.4 & \$2.50 & \$15.00 \\
\hline
\end{tabular}
\caption{Per-token pricing for each model.}
\label{tab:token-pricing}
\end{table}


\end{itemize}



\section{Project Scope}
% Even though Capstone Projects are modeled according to the same patterns as larger-scale projects, they are comparatively limited because of time constraints. To account for this, identify what is “in scope” vs. “out of scope” for the system(s) you will build (first prototype, incremental improvement, etc.) and requirements prioritization (don’t forget about data sources, scaling, etc. over time).
\subsection{In-Scope}

\subsubsection{Prototype Pipeline}
The first prototype replaces the existing stateless pipeline with a LangGraph-based agentic tool-use loop, targeting a top-5 ranking on the Spider 2.0-Lite SQLite track. Instead of a rigid multi-stage sequence, a single LLM conversation orchestrates schema exploration, sub-query testing, and self-correction iteratively. The agent is equipped with tools to inspect database metadata (tables, columns, sample rows), execute exploratory sub-queries against the SQLite file, validate intermediate results before composing the final query, and self-correct within the same conversation while preserving reasoning history.

\subsubsection{Context Engine}
A one-time offline schema enrichment step uses Qwen3-32B-AWQ to generate natural language descriptions for every table and column across the 29 benchmark databases. The enriched metadata is embedded using nomic-embed-text-v1.5 and stored in a DuckDB vector index (\texttt{dce.duckdb}). At query time, the agent calls a \texttt{search\_context} tool to retrieve the most relevant schema chunks via cosine similarity, reducing reliance on full schema dumps in the prompt.

\subsubsection{Incremental Improvements}
Building on the prototype, planned improvements include integrating the existing primary and foreign key module for join path injection, query expansion with consensus voting for schema retrieval, SQLite-specific prompt guardrails derived from failure analysis on benchmark runs, and a multi-subagent framework where specialized agents handle query planning, SQL generation, and validation in parallel before a merger agent composes the final query.


\subsubsection{Evaluation}
Evaluation targets the Spider 2.0-Lite SQLite track (135 questions across 29 databases) using Execution Accuracy as the primary metric. We benchmark against the current 38\% baseline, with a target of exceeding 55\%. Model progression moves from local open-source models (Qwen3-8B, Qwen3-32B) through to frontier APIs (Gemini Pro, GPT-4o) as the system matures.

\subsection{Out-of-Scope}

\subsubsection{Production Readiness}
The framework is designed purely to maximize benchmark accuracy. Usability, latency optimization, cost efficiency, and production deployment are not considerations for this project.

\subsubsection{Generalization}
The system is tuned specifically for Spider 2.0-Lite's database sizes and question complexity. We do not test scaling to larger datasets, nor do we address generalizability to other SQL dialects (BigQuery, Snowflake) or other benchmarks beyond Spider 2.0-Lite SQLite.

% --------------------
\section{Terminology, Definitions, Acronyms, and Abbreviations}
% --------------------
%Include all definitions, acronyms, and abbreviations necessary to understand your solution easily. You can use a table to organize your definitions if necessary. For example, define the NFRs, the metrics that are going to be used in the project, words associated with Data Structures (e.g., what is a feature vector?), terms used in the project (e.g., what is meant by a classifier?), etc.

\begin{itemize}

\item \textbf{SQL}: Structured Query Language; the standard programming language for managing and querying relational databases.
\item \textbf{NL / NLP}: Natural Language / Natural Language Processing; refers to human-spoken or written language and the computational methods used to interpret it.
\item \textbf{NL2SQL}: An abbreviation of Natural Language to SQL.
\item \textbf{LLM}: Large Language Model; a type of AI (e.g., Gemini 3, GPT-4) trained on vast datasets to understand and generate text.
\item \textbf{Agentic System}: An AI architecture where models act as “agents” that can reason, use tools, execute code, and self-correct rather than providing a single-shot response.
\item \textbf{Spider 2.0}: A rigorous benchmark for Text-to-SQL evaluation focusing on real-world enterprise complexity, including large schemas and multiple dialects. ``Lite'' refers to the subset of the benchmark used for faster iteration, while ``Snow'' specifically targets Snowflake-specific challenges.
\item \textbf{Execution Accuracy (EX)}: The primary metric for success, measuring whether the predicted SQL query, when executed, returns the exact same result set as the ground truth query.
\item \textbf{Dialect}: Variations of SQL specific to different database engines (e.g., Google BigQuery, Snowflake, SQLite) which often have unique syntax for dates, window functions, and metadata.
\item \textbf{RAG}: Retrieval-Augmented Generation; a technique where the system retrieves relevant documents or schema metadata from a vector database to provide context to the LLM.
\item \textbf{Schema}: The structural blueprint of a database, including table names, column names, data types, and the relationships (keys) between them.
\item \textbf{Vector Embedding}: A numerical representation of text (like a column description) that allows the system to find ``semantically similar'' items using mathematical distance.
\item \textbf{Self-Correction}: An iterative process where an agent examines the error logs or output of a generated query and modifies the code to fix syntax or logic flaws.
    \item \textbf{LangGraph}: An open source framework to build, manage, and deploy agentic AI workflows.
    \item \textbf{Knowledge Graph}: A graph-structured representation of a dtabase schema
    \item \textbf{Context Engine}: A hybrid retrieval system that stores table and column chunks as vector embeddings.
    \item \textbf{CTE}: A temporary, named result set in SQL.
    \item \textbf{Query planning}: A step where the user's natural language is broken down into subtasks for a final SQL generation
    \item \textbf{Schema linking}: A process to map natural language terms to database objects.
\end{itemize}

\section{Reflection}
\begin{itemize}
    \item \textbf{Do you anticipate these requirements would change as you do more research and get more data?} \\
    Yes, several requirements are likely to evolve as we gather empirical data from actual pipeline runs. The accuracy target is the most dynamic of these. Early benchmark runs with Qwen3-8B on the Spider 2.0-Lite SQLite track are already reaching around 40\%, which is encouraging for a smaller local model. The final accuracy we can achieve will depend heavily on our experiments with larger local models (Qwen3-32B, Qwen3-72B) and eventually frontier APIs, so the 55\% target should be seen as a floor rather than a ceiling. Cost estimates will also shift as we learn how many self-correction iterations the agent takes per question on average and how that scales across a full 135-question run.

    \item \textbf{How did you differentiate between the functional and non-functional requirements and decide which ones to select for your project?} \\
    We used a straightforward test: if a requirement described something the system must do for a user (a capability or behavior), it was functional; if it described a constraint on how well the system does something, it was non-functional. Functional requirements like SQL generation, schema retrieval, and self-correction describe discrete, testable behaviors. Non-functional requirements like latency, scalability, and transparency describe performance envelopes that cut across multiple functional areas. Security was a borderline case because it could be framed as a feature (a validation step that rejects write queries) or as a system-wide constraint. We treated it as non-functional because it applies uniformly to every SQL generation flow rather than being a standalone deliverable.

    \item \textbf{How do intended users affect your requirements, and if the users change, do you think that you will need to update the requirements?} \\
    The two user groups shaped the requirements in concrete ways. ML Researchers motivated the modularity requirement (components should be independently replaceable for ablation studies), the reproducibility requirement (deterministic run IDs, config hashing), and detailed logging (correction traces, error taxonomy). Without a research user, these could be simplified to a single output file per run. Data and Analytics Teams motivated the scalability requirement (handling schemas with 1000+ columns) and the accuracy target, since a system that fails on complex real-world schemas would not serve business users. If we added a third user group with production needs, such as availability, concurrent user support, or API rate limiting, we would need a significant set of new non-functional requirements that are explicitly out of scope today.

    \item \textbf{Have you taken broader reach and accessibility into consideration during the design process?} \\
    Our design is primarily optimized for benchmark evaluation, which means usability and accessibility are acknowledged trade-offs rather than oversights. That said, some design choices do move in the right direction: human-readable trace logs make the system's reasoning inspectable for non-technical stakeholders, and the confidence indicator (HIGH / MEDIUM / LOW) on returned results gives users a signal about whether to trust the output without needing to understand the underlying SQL. Multilingual support and natural language error explanations would meaningfully extend accessibility for Data and Analytics Teams, but these fall outside the current project scope and would require a separate user research effort to design well.

    \item \textbf{What would you do differently next time?} \\
    We would start with system context and data flow diagrams before writing any prose requirements. This time, we wrote pipeline stage descriptions first and then revised them multiple times when diagrams exposed gaps, such as where semantic validation fits relative to execution, or how early stopping interacts with the correction loop. Starting visual made these dependencies explicit immediately. We would also instrument the pipeline from the very first run to capture per-stage latency and token usage, which would have made the latency and cost requirements much more grounded from the start rather than estimated after the fact.
\end{itemize}

\section{Appendix}
% This section contains any additional information you’d like to preserve in this document for context. For example, consider having a Glossary, or any additional materials you discovered or created in the process of making this document.
\subsection{Spider 2.0 Dataset and Leaderboard}

\begin{description}
    \item[Spider 2.0 Benchmark] A large-scale, enterprise-level benchmark for 
    text-to-SQL translation. Unlike academic predecessors, it contains 687 databases 
    and over 1{,}000 tables, designed to simulate professional data environments with 
    complex schemas and cross-database reasoning.
    \item[Spider 2.0 Lite] A curated subset (approximately 547 tasks) that maintains 
    high complexity while allowing for more rapid and cost-effective iteration during 
    development.
    \item[Spider 2.0 Leaderboard] A public ranking system that tracks the performance 
    of models and agentic frameworks (e.g., ReFoRCE \cite{deng2025reforce}, DAIL-SQL). It serves as the 
    primary standard for state-of-the-art performance in the field.
    \item[Dialects (BigQuery, Snowflake, SQLite)] The three SQL environments supported 
    by Spider 2.0. Each requires different handling of metadata, data types (e.g., JSON 
    flattening in Snowflake), and specific syntax (e.g., backticks in BigQuery).
\end{description}

\subsection{Sample Execution Log}

A brief example of a first-pass failure and a successful second-pass correction, 
illustrating the self-correction loop described in Section~2:

\begin{quote}
\textbf{Question:} ``Find the average revenue for 2023.'' \\[0.5em]
\textbf{Iteration 0 (Error):} \texttt{SELECT AVG(rev) FROM sales WHERE year = 2023} 
$\rightarrow$ \textit{``Column `rev' does not exist.''} \\[0.3em]
\textbf{Correction Agent:} Analyzes schema $\rightarrow$ Finds 
\texttt{total\_revenue} column. \\[0.3em]
\textbf{Iteration 1 (Success):} \texttt{SELECT AVG(total\_revenue) FROM sales 
WHERE year = 2023} $\rightarrow$ \textit{Execution Success.}
\end{quote}

\subsection{Pipeline Configuration Parameters}

\begin{description}
    \item[\texttt{max\_corrections} (default: 3)] Maximum self-correction iterations 
    per query.
    \item[\texttt{query\_rewriting} (default: False)] Enable question rewriting before 
    schema retrieval.
    \item[\texttt{query\_planning} (default: False)] Enable CTE/join strategy planning 
    before generation.
    \item[\texttt{semantic\_validation} (default: False)] Enable pre/post-execution 
    semantic validation.
    \item[\texttt{enhanced\_correction} (default: False)] Use error classification and 
    progressive guidance in the correction loop.
    \item[\texttt{early\_stopping} (default: False)] Skip correction for 
    already-successful queries in parallel mode.
    \item[\texttt{num\_parsers} (default: 3)] Number of independent parser agents in 
    column retrieval voting.
    \item[\texttt{max\_workers} (default: 8)] Thread pool size for parallel execution.
\end{description}

\subsection{Error Taxonomy}

The enhanced correction module classifies execution failures into the following 
categories, used to select error-specific correction strategies:

\begin{description}
    \item[\texttt{SYNTAX\_ERROR}] Malformed SQL (missing keywords, unmatched 
    parentheses).
    \item[\texttt{SCHEMA\_ERROR}] References to non-existent tables or columns.
    \item[\texttt{AMBIGUOUS\_REFERENCE}] Unqualified column names that exist in 
    multiple joined tables.
    \item[\texttt{TYPE\_ERROR}] Data type mismatches (e.g., comparing string to 
    integer).
    \item[\texttt{LOGICAL\_ERROR}] Query executes but returns semantically incorrect 
    results.
    \item[\texttt{CONSTRAINT\_ERROR}] Violations of database constraints (e.g., 
    NOT NULL, UNIQUE).
    \item[\texttt{PERFORMANCE\_ERROR}] Query exceeds execution timeout.
    \item[\texttt{GENERAL\_ERROR}] Catch-all for unclassified failures.
\end{description}
\section{Changes To Previous Deliverables}
% Use this section to outline any changes you had to make in your previous documents, the Vision Document, because of your activities as part of this deliverable.
No changes were made in the Vision Document as a result of the activities in the Requirements Document. We only modified the formatting to incorporate the teaching staff feedback.  
\section{Response to Teaching Staff Feedback}
% After receiving feedback from the instructor team on this deliverable,
% please provide a response to the instructor team’s feedback on this
% deliverable, clearly outlining the revisions and improvements you have
% implemented to address their comments.

% First you should address \emph{in the text itself},  the comments that have been added to your
% report by the reviewers, making any necessary corrections, deletions
% or additions as requested or suggested.  If there is a comment that you disagree with
% it, indicates so with your justification below.


% Once you address all the comments, please list below as an enumerated
% list your responses to each comment (e.g., what you added, what you
% deleted, what you changed, why you disagree with a comment, etc.),
% identifying in some suitable way what you are responding to
% (e.g. ``Comment in Section 1 paragraph 2..'')

% The TAs will review these to verify that all comments are acceptably addressed.

The comment on Section~\ref{sec:nfr} Non-Functional Requirements was: \texttt{The remaining deduction reflects that these enforcement mechanisms are not always explicitly structured or consistently stated for each requirement, rather than an absence of content.}

\begin{enumerate}
\item Comment in Section 1 (Intended Users): There is still some room for improvement in explicitly structuring the mapping between user groups, their use cases, and corresponding system requirements. This is the main reason for the remaining deduction.
\\Response: We formatted the section to explicitly state the use case and requirements ff the user groups.
\item Comment in References Section: The section includes appropriate and correctly formatted references that align with the in-text citations. The only issue is the lack of numbering, which affects consistency with standard citation formatting.
\\Response: We fixed the document formatting to include numbered citations.
\item Comment in Section 6 (Terminology, Definitions, Acronyms, and Abbreviations): The section provides a strong and well-organized set of definitions covering most key concepts in the project, which supports clarity across the document. However, it is missing some system-specific terminology and includes minor redundancy.
\\Response: Added system-specific terminology like LangGraph, Query Planning etc. and removed redundant items.
\item Comment in Section~\ref{sec:nfr}: Added a more explicit enforcement mechanism for latency.
\item Comment in Section Resource Requirements: \textit{" There are still minor areas that could be improved (e.g., slightly more detail on how resources scale across experiments or supporting components like vector storage), but overall the section is solid and meets expectations."} \\
Expanded the Hardware item to describe how resource requirements evolve across three phases (local open-source models, larger quantised models, frontier APIs) and added details on the vector storage component (\texttt{dce.duckdb}) and its one-time preprocessing cost.
\item Comment in Section Project Scope: \textit{"That said, the section remains somewhat high-level and does not explicitly define concrete deliverables or structure the scope in more detail."} \\
Broke down both In-Scope and Out-of-Scope into concrete subsubsections (Prototype Pipeline, Context Engine, Incremental Improvements, Evaluation; and Production Readiness, Generalization) to make deliverables and exclusions more explicit.
\item Comment in Section Reflection: \textit{"Some responses are relatively brief and could be expanded with more concrete examples and deeper discussion of challenges or trade-offs."} \\
Expanded all five reflection answers with concrete examples drawn from actual project experience: cited observed latency numbers (60-100s vs. the 30s target), specific benchmark failure categories (RFM, regression), the confidence indicator as a concrete accessibility feature, and the instrumentation gap as a lesson learned.




  
\end{enumerate}
\bibliography{references}
